\cleardoublepage
\phantomsection
\pdfbookmark{Abstract}{Abstract}
\begingroup
\let\clearpage\relax
\let\cleardoublepage\relax
\chapter*{Abstract}
\noindent
This thesis investigates the use of Structured Logic Tensor Networks (LTNs) for neuro-symbolic learning.
The goal is to combine neural learning with explicit background logical knowledge.
The proposed approach addresses an AI classification task with limited and heterogeneous training data.
Entities and concepts are represented through learnable continuous embeddings.
Logical relations and domain constraints are expressed using first-order logic.
These logical formulas are converted into differentiable constraints using fuzzy logic semantics.
The model jointly learns from labeled examples and the available background knowledge.
A structured representation is used to organize concepts, relations, and logical dependencies.
This allows domain knowledge to guide the learning process rather than relying only on observed data.
The approach is evaluated against conventional neural network baselines.
Experiments investigate both predictive accuracy and logical consistency of the learned model.
Additional experiments analyze the effect of increasing or reducing the amount of background knowledge.
The results examine whether logical constraints improve generalization, particularly with limited training data.
The thesis also discusses computational and modeling challenges associated with differentiable logical reasoning. Overall, the study demonstrates the potential of Structured LTNs for integrating learning and reasoning in AI systems.
\endgroup
\vfill